\section{Methodology}
\label{sec:method}

To address the cross-project generalisation gap and the data-silo problem simultaneously, we propose \emph{VulMorph-Fed}. The framework replaces standard parameter averaging with a privacy-preserving federated aggregation of semantic structural prototypes. It consists of three components: (1)~Vulnerability-Critical Subgraph Abstraction (VCSA), which maps functions onto project-invariant structural morphologies; (2)~Morphology-Conditioned Federated Prototype Aggregation (MCFPA), which aggregates structural knowledge under differential privacy; and (3)~Morphology-Guided Message Passing (MGMP), which injects the aggregated global knowledge into local graph learning. Figure~\ref{fig:worked_example} illustrates the end-to-end abstraction on two lexically different but structurally similar vulnerable functions.

\subsection{Graph Construction}
\label{sec:graph_construction}

Let $\mathcal{G}_k = \{(G_i^{(k)}, y_i^{(k)}, c_i^{(k)})\}_{i=1}^{N_k}$ denote the local dataset of client $k \in \{1, \dots, K\}$, where $y_i \in \{0,1\}$ is the binary vulnerability label and $c_i$ the CWE annotation of vulnerable functions (Section~\ref{sec:label_model}).

Full Code Property Graph (CPG) extraction requires a heavyweight static-analysis engine such as Joern~\cite{yamaguchi2014modeling}, which is impractical to deploy at every federated edge client and complicates exact reproduction. Our implementation instead parses every function with \texttt{tree-sitter} --- a production-grade, error-tolerant C/C++ parser --- and builds the graph $G = (\mathcal{V}, \mathcal{E})$ directly from the abstract syntax tree:
\begin{itemize}
    \item one node per \emph{named AST node}, in pre-order, capped at $|\mathcal{V}| \le 200$;
    \item \emph{syntax edges} $\mathcal{E}_{AST}$: parent--child links in both directions;
    \item \emph{order edges} $\mathcal{E}_{SIB}$: consecutive named siblings (source-order sequencing, the AST analogue of Devign's NCS edges~\cite{zhou2019devign});
    \item \emph{data-dependence proxy} edges $\mathcal{E}_{DD}$: def-use chains linking successive occurrences of the same identifier leaf.
\end{itemize}
The edge set is $\mathcal{E} = \mathcal{E}_{AST} \cup \mathcal{E}_{SIB} \cup \mathcal{E}_{DD}$; on our corpora this yields on average $\sim$130 nodes and $\sim$350 edges per function. Functions that fail to parse (a small minority) fall back to a token-level sequence graph and are marked as such. These graphs carry real syntactic structure, though not the interprocedural control- and data-flow precision of a compiler-grade CPG; Section~\ref{sec:discussion} discusses this trade-off, and graph construction is isolated behind a single interface so Joern-derived CPGs can be substituted without touching the rest of the framework.

\subsection{Vulnerability-Critical Subgraph Abstraction (VCSA)}
\label{sec:vcsa}

Raw code graphs carry thousands of project-specific API and identifier tokens, causing GNNs to overfit local lexical distributions. VCSA replaces the lexical node labels with a small taxonomy of abstract operation types and learns to down-weight vulnerability-irrelevant edges.

\subsubsection{Morphological taxonomy}
\label{sec:taxonomy}

We define a fixed semantic taxonomy $\mathcal{T}$ of $|\mathcal{T}| = 8$ operation types, listed exhaustively in Table~\ref{tab:taxonomy} together with the reserved fallback type \texttt{UNKNOWN}. A node therefore takes one of $|\mathcal{T}| + 1 = 9$ morphological categories --- the \texttt{UNKNOWN} type has its own feature slot. For the taxonomy-size sensitivity analysis (RQ2b) we additionally define strictly nested refinements with $|\mathcal{T}| = 16$ and $|\mathcal{T}| = 32$ (e.g., \texttt{MEMORY\_ACCESS} splits into \texttt{MEMORY\_ALLOC}, \texttt{MEMORY\_FREE}, \texttt{MEMORY\_COPY}, \texttt{STRING\_OP}); the complete lists ship with the source code.

Each node's feature vector is the sum of two embeddings, both drawn from \emph{project-invariant} vocabularies: (i) its morphological type in $\mathcal{T} \cup \{\texttt{UNKNOWN}\}$, and (ii) its \emph{grammar kind} --- the tree-sitter node category (\texttt{call\_expression}, \texttt{subscript\_expression}, \texttt{if\_statement}, \dots), a fixed set of 364 categories determined entirely by the C grammar and therefore identical for every project and client by construction. No identifier, API name, literal or any other project-specific token ever enters the network.

\begin{table}[t]
\centering
\caption{The complete $|\mathcal{T}| = 8$ morphological taxonomy (plus the reserved \texttt{UNKNOWN} fallback). A node takes one of $|\mathcal{T}|+1 = 9$ morphological \emph{categories}; each category is mapped to a learned embedding (dimension 64 in our configuration), summed with the node's grammar-kind embedding.}
\label{tab:taxonomy}
\resizebox{\linewidth}{!}{%
\begin{tabular}{ll}
\toprule
Abstract type & Covers \\
\midrule
\texttt{MEMORY\_ACCESS} & allocation/free/copy/set and string APIs (\texttt{malloc}, \texttt{free}, \texttt{memcpy}, \texttt{strcpy}, \dots) \\
\texttt{ARRAY\_INDEX}   & array subscript operations \texttt{a[i]} \\
\texttt{PTR\_DEREF}     & pointer dereference and member access (\texttt{*p}, \texttt{p->f}, unary \texttt{\&}) \\
\texttt{CONTROL\_BRANCH}& \texttt{if}/\texttt{else}/\texttt{switch}, loops, \texttt{goto}/\texttt{break}/\texttt{continue}/\texttt{return}, ternary \\
\texttt{ARITH\_OP}      & arithmetic and bitwise operators (\texttt{+ - * / \% \& | \^{} \~{} << >>}) \\
\texttt{COMPARISON}     & relational and logical operators (\texttt{== != < > <= >= \&\& || !}) \\
\texttt{CALL\_SITE}     & any other function/macro call (user-defined and unresolved calls included) \\
\texttt{ASSIGN}         & plain and compound assignments (\texttt{=}, \texttt{+=}, \dots) \\
\midrule
\texttt{UNKNOWN}        & all remaining tokens (identifiers, literals, punctuation, declarations) \\
\bottomrule
\end{tabular}}
\end{table}

\subsubsection{The abstraction function $\phi_{\mathrm{type}}$}

The mapping $\phi_{\mathrm{type}}: \mathcal{V} \rightarrow \mathcal{T} \cup \{\texttt{UNKNOWN}\}$ is a deterministic rule set evaluated on the AST, in priority order:
\begin{enumerate}
    \item a \texttt{call\_expression} is classified by its callee name in two stages: an exact match against fixed libc allow-lists (\texttt{malloc}, \texttt{memcpy}, \texttt{strcpy}, \dots), then a set of \emph{affix rules} that recognise project-wrapped equivalents. The affix stage is essential rather than cosmetic: production C code almost never calls libc allocators directly (FFmpeg uses \texttt{av\_malloc}, GLib \texttt{g\_malloc}, the Linux kernel \texttt{kmalloc}, OpenSSL \texttt{OPENSSL\_malloc}), so an exact-match list alone routes nearly every real allocation site to \texttt{CALL\_SITE}. We measured this: on Devign, exact matching alone assigns \texttt{MEMORY\_ACCESS} to only 0.23\% of AST nodes, rising to 0.72\% once affix rules are applied --- i.e.\ two thirds of the memory-safety signal the taxonomy is designed to capture is invisible without them. The rules test for the affixes \texttt{*alloc*}, \texttt{*free*}, \texttt{*memcpy*}, \texttt{*memset*}, \texttt{*str\{cpy,cat,len\}*} and are evaluated in an order that resolves \texttt{realloc} and \texttt{dealloc} before the generic \texttt{alloc} test. Only the matched \emph{class} is retained --- never the identifier --- so project invariance is preserved. Unmatched calls, including all user-defined functions, macros and unresolved calls, map to \texttt{CALL\_SITE}; no name resolution or build system is required;
    \item statement and expression kinds map directly: \texttt{subscript\_expression} $\to$ \texttt{ARRAY\_INDEX}; \texttt{field\_expression} and unary pointer expressions $\to$ \texttt{PTR\_DEREF}; \texttt{if}/\texttt{switch}/loop/jump statements and ternaries $\to$ \texttt{CONTROL\_BRANCH}; \texttt{assignment\_expression} and initialising declarators $\to$ \texttt{ASSIGN};
    \item a \texttt{binary\_expression} maps by its operator text to \texttt{ARITH\_OP} or \texttt{COMPARISON} (Table~\ref{tab:taxonomy});
    \item three further kind rules complete the set: \texttt{cast\_expression} $\to$ \texttt{PTR\_DEREF} (a cast reinterprets a pointee type), \texttt{update\_expression} (\texttt{i++}, \texttt{i--}) $\to$ \texttt{ARITH\_OP}, and unary \texttt{!} $\to$ \texttt{COMPARISON} with the remaining unary operators $\to$ \texttt{ARITH\_OP};
    \item every remaining node maps to \texttt{UNKNOWN}.
\end{enumerate}
An equivalent token-context rule set (documented in the code) covers the rare parse-failure fallback graphs. The rules are implemented once at the finest ($|\mathcal{T}|=32$) granularity and projected onto the coarser taxonomies, so all three taxonomy sizes share a single rule set. $\phi_{\mathrm{type}}$ is a discrete label assignment and is \emph{not} differentiable; the only learned part of VCSA is the edge-weighting module below.

On our four evaluation corpora, 25--26\% of AST nodes receive a non-\texttt{UNKNOWN} morphological type (the remainder are declarations, identifiers, literals and other syntax carriers, which still contribute their grammar-kind feature). The abstraction does not delete nodes: all nodes are retained with their graph structure, but their features are collapsed from a $\sim$10,000-token lexical vocabulary onto $|\mathcal{T}|+1$ morphological categories plus 364 grammar kinds. This many-to-one collapse is what removes project-specific lexical information from everything the model computes and, downstream, from everything a client transmits (Section~\ref{sec:privacy}).

\subsubsection{Learned edge weighting}

VCSA additionally learns a soft edge-importance mask. A two-layer MLP over concatenated endpoint features produces
\begin{equation}
    \varepsilon_e = \sigma\!\left(\mathrm{MLP}([\mathbf{x}_u \,\|\, \mathbf{x}_v])\right) \in [0,1] \quad \text{for } e = (u,v) \in \mathcal{E},
\label{eq:edge_mask}
\end{equation}
which reweights messages during propagation (Eq.~\ref{eq:mgmp_local}); edges are softly down-weighted rather than hard-pruned, keeping the pipeline trainable end-to-end. Sparsity is encouraged by an $\ell_1$ penalty $\mathcal{L}_{\mathrm{sparse}} = \frac{1}{|\mathcal{E}|}\sum_e \varepsilon_e$. To align morphologies across clients we further use a structural contrastive loss over graph embeddings $\mathbf{h}_G$:
\begin{equation}
    \mathcal{L}_{SCL} = \frac{1}{|\mathcal{P}^+ \cup \mathcal{P}^-|} \left[ \sum_{(i,j) \in \mathcal{P}^+}\!\! \|\mathbf{h}_{G_i} - \mathbf{h}_{G_j}\|_2^2
    + \!\!\sum_{(i,j) \in \mathcal{P}^-}\!\! \max\left(0, m - \|\mathbf{h}_{G_i} - \mathbf{h}_{G_j}\|_2^2\right) \right],
\label{eq:scl}
\end{equation}
where $\mathcal{P}^+$ are same-CWE vulnerable pairs and $\mathcal{P}^-$ are (vulnerable, benign) pairs within a mini-batch, with margin $m = 1$. The total local objective is $\mathcal{L} = \mathcal{L}_{BCE} + \alpha \mathcal{L}_{SCL} + \gamma \mathcal{L}_{\mathrm{sparse}}$ with $\alpha = 0.1$, $\gamma = 0.01$; $\mathcal{L}_{BCE}$ is class-weighted binary cross-entropy with $\mathrm{pos\_weight} = \min(N^-/N^+, 20)$ computed on each client's local split.

\begin{figure}[t]
\centering
\begin{minipage}{0.48\linewidth}
\footnotesize
\textbf{Project A (FFmpeg-style):}
\begin{verbatim}
n = av_get_size(ctx);
buf = av_malloc(n);
for (i = 0; i <= n; i++)
    buf[i] = src[i];
\end{verbatim}
\end{minipage}\hfill
\begin{minipage}{0.48\linewidth}
\footnotesize
\textbf{Project B (OpenSSL-style):}
\begin{verbatim}
len = BN_num_bytes(a);
out = OPENSSL_malloc(len);
for (k = 0; k <= len; k++)
    out[k] = in[k];
\end{verbatim}
\end{minipage}
\vspace{2pt}
\begin{minipage}{\linewidth}
\centering
\footnotesize
\textbf{Morphology sequence emitted by $\phi_{\mathrm{type}}$ --- identical for both fragments:}\\[2pt]
\texttt{CONTROL\_BRANCH} $\rightarrow$ \texttt{ASSIGN} $\rightarrow$ \texttt{ASSIGN} $\rightarrow$ \texttt{ASSIGN} $\rightarrow$ \texttt{COMPARISON} $\rightarrow$ \texttt{ARITH\_OP} $\rightarrow$ \texttt{CALL\_SITE} $\rightarrow$ \texttt{MEMORY\_ACCESS} $\rightarrow$ \texttt{ASSIGN} $\rightarrow$ \texttt{ARRAY\_INDEX} $\rightarrow$ \texttt{ARRAY\_INDEX}
\end{minipage}
\caption{Worked example of the VCSA abstraction. Two off-by-one buffer overflows from different projects share no API or identifier token --- \texttt{av\_get\_size}/\texttt{BN\_num\_bytes}, \texttt{av\_malloc}/\texttt{OPENSSL\_malloc}, \texttt{buf}/\texttt{out} --- yet $\phi_{\mathrm{type}}$ maps both onto a byte-identical morphology: a size query (\texttt{CALL\_SITE}), an allocation (\texttt{MEMORY\_ACCESS}), and a loop whose \emph{inclusive} bound (\texttt{COMPARISON} on $\le$) drives an array write (\texttt{ARRAY\_INDEX}). The sequence shown is the literal output of the released \texttt{classify\_token}/\texttt{\_classify\_ast\_node} implementation on these two inputs (breadth-first over the AST, \texttt{UNKNOWN} nodes elided), not an idealisation; it is regenerated from the code so that figure and implementation cannot diverge. Note that recognising \texttt{av\_malloc} and \texttt{OPENSSL\_malloc} as allocations requires the affix rules of Section~\ref{sec:vcsa} --- an exact libc allow-list maps both to \texttt{CALL\_SITE} and the shared morphology is lost.}
\label{fig:worked_example}
\end{figure}

\subsection{Label Model}
\label{sec:label_model}

Real vulnerability corpora have a many-to-many CVE--function--CWE structure, so we state our label model explicitly:
\begin{itemize}
    \item \textbf{Detection is strictly binary.} The classification head predicts $\hat{y} \in [0,1]$ (vulnerable vs.\ benign); all reported Precision, Recall, F1 and AUC values are computed from this binary head. The decision threshold is calibrated on held-out training-project data as described in Section~\ref{sec:experiment}, never on the test set; AUC is threshold-free. CWE information never enters the detection head.
    \item \textbf{CWE labels condition only the prototypes.} A vulnerable function annotated with several CWEs contributes to the prototype of its \emph{first-listed (primary)} CWE. Benign functions contribute to a dedicated \emph{benign} prototype slot, giving the bank an explicit model of non-vulnerable morphology.
    \item \textbf{CWE vocabulary.} The prototype bank has $|\mathcal{C}| + 1 = 11$ slots: the nine most frequent CWE types in the training corpus receive dedicated slots, all remaining types share an \texttt{OTHER} slot, and the eleventh slot is the benign prototype. This bucketing is computed once per dataset and reported with the released configurations. Datasets without per-function CWE annotations (Devign) degenerate gracefully to a single vulnerable-prototype slot plus the benign slot.
\end{itemize}

\subsection{Morphology-Conditioned Federated Prototype Aggregation (MCFPA)}
\label{sec:mcfpa}

FedAvg-style parameter averaging~\cite{mcmahan2017fedavg} assumes near-IID clients; vulnerability data is extremely non-IID in its CWE composition. MCFPA instead federates \emph{class-conditioned prototypes} over the $|\mathcal{C}|+1$ bank slots of Section~\ref{sec:label_model} (CWE buckets plus the benign slot). After local training, client $k$ computes, for each slot $c$ with local support $N_{c,k} > 0$,
\begin{equation}
    \mathbf{p}_c^{(k)} = \frac{1}{N_{c,k}} \sum_{\substack{G \in \mathcal{G}_k:\; y_G = 1,\; c_G = c}} \mathrm{clip}_R\!\left(f_{\theta_k}(\tilde{G})\right),
    \qquad
    \mathrm{clip}_R(\mathbf{z}) = \mathbf{z} \cdot \min\!\left(1, \tfrac{R}{\|\mathbf{z}\|_1}\right),
\label{eq:prototype}
\end{equation}
i.e., the centroid of the $\ell_1$-clipped embeddings of its local functions assigned to slot $c$ (vulnerable functions of the corresponding CWE bucket; all benign functions for the benign slot). Slots with $N_{c,k} = 0$ are transmitted as zero rows and carry no information.

\subsubsection{Differential privacy}
\label{sec:privacy}

Before upload, each non-empty prototype row is perturbed with the Laplace mechanism:
\begin{equation}
    \hat{\mathbf{p}}_c^{(k)} = \mathbf{p}_c^{(k)} + \mathrm{Lap}\!\left(0,\; \frac{2R}{N_{c,k}\,\varepsilon}\right)^{\otimes d}.
\label{eq:laplace}
\end{equation}

VulMorph-Fed applies two mechanisms, so that the guarantee covers the whole
pipeline rather than only the released artefact.

\paragraph{Mechanism 1: private encoder training.}
Local training uses DP-SGD~\cite{abadi2016deep}: each per-example gradient is
clipped to $\|g_i\|_2 \le C$ and the batch sum is perturbed with Gaussian noise
of standard deviation $\sigma C$ before the optimiser step. The encoder
$\theta_k$ is therefore itself a differentially private function of the
client's data. This matters: without it, everything below would be conditional
on an encoder that had seen the raw data unprotected, and changing one record
could move every embedding by an unbounded amount.

\paragraph{Mechanism 2: private prototype release.}
Prototypes are built from $\ell_1$-clipped embeddings (Eq.~\ref{eq:prototype})
and perturbed per class (Eq.~\ref{eq:laplace}).

\begin{proposition}[Per-round privacy of the prototype release]
\label{prop:dp}
Fix a client $k$ and condition on its (DP-trained) encoder. Let two
neighbouring datasets differ in one function. Releasing the noisy bank
$\{\hat{\mathbf{p}}_c^{(k)}\}$ of Eq.~\eqref{eq:laplace} satisfies
$2\varepsilon$-differential privacy with respect to that record.
\end{proposition}

\begin{proof}
With the encoder fixed, each per-sample embedding lies in the $\ell_1$ ball of
radius $R$, so replacing one record changes an affected class mean by at most
$\Delta_1 = 2R/N_{c,k}$ in $\ell_1$ norm; the Laplace mechanism with scale
$\Delta_1/\varepsilon$ is $\varepsilon$-DP for that row~\cite{dwork2014algorithmic}.
Under the replacement neighbouring relation a record change can affect
\emph{two} rows --- it leaves the slot of its old label and joins the slot of
its new one --- so parallel composition across slots does not apply and the two
affected rows compose sequentially, giving $2\varepsilon$. (A restricted
neighbouring relation that preserves the label would give $\varepsilon$; we
report the weaker, safer constant.)
\end{proof}

\begin{theorem}[End-to-end budget]
\label{thm:total}
Let $\varepsilon_{\mathrm{sgd}}$ be the budget consumed by DP-SGD over all local
steps, accounted under R\'enyi differential privacy for the subsampled Gaussian
mechanism~\cite{mironov2019rdp} at a target $\delta$. Then $T$ rounds of
VulMorph-Fed satisfy
\begin{equation}
    \varepsilon_{\mathrm{tot}} = \varepsilon_{\mathrm{sgd}} + 2T\varepsilon
\label{eq:total_eps}
\end{equation}
$(\varepsilon_{\mathrm{tot}}, \delta)$-differential privacy with respect to any
single function in a client's dataset, covering both the transmitted prototypes
and the locally trained encoder.
\end{theorem}

\begin{proof}
DP-SGD over the local steps is $(\varepsilon_{\mathrm{sgd}}, \delta)$-DP by the
R\'enyi accountant; the $T$ prototype releases are $2T\varepsilon$-DP by
Proposition~\ref{prop:dp} and sequential composition. The two mechanisms
observe the same private dataset, so their budgets add.
\end{proof}

We use the standard R\'enyi accountant (the binomial bound of
Mironov et al.~\cite{mironov2019rdp} for the sampled Gaussian, converted to
$(\varepsilon,\delta)$ over a grid of orders) rather than advanced composition,
which would overstate $\varepsilon_{\mathrm{sgd}}$ by more than an order of
magnitude; our implementation reproduces the reference regime of
Abadi et al.~\cite{abadi2016deep} to within the expected margin.

\paragraph{Residual leakage outside the analysis.}
Two channels remain, and we state them rather than leaving them implicit. The
noise scale $2R/(N_{c,k}\varepsilon)$ depends on the private count $N_{c,k}$,
and the \emph{support pattern} of a client's occupied slots is transmitted
unnoised, revealing which weakness classes a client possesses at all. Padding
the bank to a fixed support and using a count-independent scale would close
both at some utility cost; we quantify neither here.

We emphasise why the morphological abstraction matters here: because the encoder consumes only $|\mathcal{T}|+1$ abstract node types, the embeddings being clipped and averaged contain no project-specific lexical information to begin with. DP protects \emph{membership} of individual functions; the abstraction removes \emph{content} (identifiers, API names) categorically rather than statistically.

\subsubsection{Affinity-weighted aggregation}

The server aggregates the noisy banks per CWE bucket. Let $A_k \subseteq \{1..K\}$ be the clients with non-zero rows for bucket $c$. With $\bar{\mathbf{p}}_j = \hat{\mathbf{p}}_c^{(j)} / \|\hat{\mathbf{p}}_c^{(j)}\|_2$, the server computes cosine-affinity weights and the global prototype
\begin{equation}
    w_c^{(k)} = \max\!\left(0,\; \frac{1}{|A_k| - 1} \sum_{j \in A_k,\, j \neq k} \bar{\mathbf{p}}_j^\top \bar{\mathbf{p}}_k \right),
    \qquad
    \mathbf{p}_c^{\mathrm{glob}} = \frac{\sum_{k \in A_k} w_c^{(k)}\, \hat{\mathbf{p}}_c^{(k)}}{\sum_{k \in A_k} w_c^{(k)}},
\label{eq:mcfpa}
\end{equation}
falling back to the uniform mean when all affinities vanish, and to the previous round's value when no client has data for $c$. Clients whose local view of CWE $c$ agrees with the consensus contribute more; outlier (or heavily noised) views are attenuated. The ablation ``w/o MCFPA'' replaces Eq.~\eqref{eq:mcfpa} with the uniform mean.

\subsection{Morphology-Guided Message Passing (MGMP)}
\label{sec:mgmp}

Each client integrates the broadcast bank $\mathcal{P}^{\mathrm{glob}} = \{\mathbf{p}_c^{\mathrm{glob}}\}$ into its GNN. Layer $l$ first performs edge-weighted local aggregation (GCN-style normalisation $\hat{d}$ with the VCSA mask $\varepsilon$ as edge weight):
\begin{equation}
    \mathbf{h}_v^{\mathrm{loc}} = \sum_{u \in \mathcal{N}(v) \cup \{v\}} \frac{\varepsilon_{uv}}{\sqrt{\hat{d}_u \hat{d}_v}}\, \mathbf{W}^{(l)} \mathbf{h}_u^{(l)},
\label{eq:mgmp_local}
\end{equation}
then attends over the global prototypes with scaled dot-product attention:
\begin{equation}
    \mathbf{m}_v = \sum_{c \in \mathcal{C}} \mathrm{softmax}_c\!\left(\frac{(\mathbf{W}^{(l)}\mathbf{h}_v^{(l)})^\top \mathbf{p}_c^{\mathrm{glob}}}{\sqrt{d}}\right) \mathbf{p}_c^{\mathrm{glob}},
\label{eq:protoattn}
\end{equation}
and fuses the two signals through a per-node morphology gate $\beta_v = \sigma(\mathbf{w}_\beta^\top [\mathbf{W}^{(l)}\mathbf{h}_v^{(l)} \,\|\, \mathbf{x}_v^{\mathrm{morph}}])$ and a learnable per-layer balance $\lambda \in [0,1]$:
\begin{equation}
    \mathbf{h}_v^{(l+1)} = \mathrm{ReLU}\!\left((1-\lambda)\, \mathbf{h}_v^{\mathrm{loc}} + \lambda\, \beta_v\, \mathbf{m}_v\right).
\label{eq:mgmp_fuse}
\end{equation}
Nodes whose abstract type is implicated in known vulnerability morphologies receive a higher gate $\beta_v$ and thus a stronger global signal. Even a client that has never observed CWE $c$ locally can recognise its structural morphology by attending to $\mathbf{p}_c^{\mathrm{glob}}$.

The global bank additionally reaches the decision through two further channels. First, the graph-level readout is $[\mathbf{h}_G^{\mathrm{mean}} \,\|\, \mathbf{h}_G^{\mathrm{max}} \,\|\, \mathrm{cos}(\mathbf{h}_G^{\mathrm{mean}}, \mathbf{p}_c^{\mathrm{glob}})_{c \in \mathcal{C} \cup \{\mathrm{ben}\}}]$: the classifier directly sees how similar the current function is to every global morphology, including the benign one. Second, a prototype-alignment regulariser
\begin{equation}
    \mathcal{L}_{\mathrm{proto}} = \frac{1}{|B|} \sum_{i \in B} \big\| \mathbf{h}_{G_i} - \mathbf{p}_{s(i)}^{\mathrm{glob}} \big\|_2^2,
\label{eq:proto_align}
\end{equation}
where $s(i)$ is sample $i$'s bank slot (its CWE bucket, or the benign slot), pulls each client's embedding space toward the shared global geometry --- the mechanism FedProto~\cite{tan2022fedproto} showed to be essential for prototype-based FL under non-IID clients. The total local objective becomes $\mathcal{L} = \mathcal{L}_{BCE} + \alpha \mathcal{L}_{SCL} + \gamma \mathcal{L}_{\mathrm{sparse}} + \mu \mathcal{L}_{\mathrm{proto}}$ with $\mu = 0.5$.

\subsection{Training and Inference Protocol}
\label{sec:protocol}

Algorithm~\ref{alg:vulmorph_fed} summarises one federation. Note that \emph{model parameters are never aggregated}: each client keeps its own encoder $\theta_k$, and the only shared state is the prototype bank.

\textbf{Inference.} The deployed global detector is the uniform probability ensemble of the $K$ client models, each conditioned on the same final prototype bank:
\begin{equation}
    \hat{p}(y \mid G) = \frac{1}{K} \sum_{k=1}^{K} \sigma\!\left(f_{\theta_k}(\tilde{G};\, \mathcal{P}^{\mathrm{glob}})\right).
\label{eq:ensemble}
\end{equation}
All headline metrics are computed from Eq.~\eqref{eq:ensemble} on the held-out cross-project test set; we additionally report the mean and standard deviation of per-client F1 to expose client-level variance. In a production deployment, Eq.~\eqref{eq:ensemble} corresponds to clients exposing scoring endpoints, or --- when a single model is required --- any client model can be used alone, at the cost quantified by the reported per-client variance.

\begin{algorithm}[t]
\caption{VulMorph-Fed}
\label{alg:vulmorph_fed}
\begin{algorithmic}[1]
\REQUIRE $K$ clients, rounds $T$, per-round budget $\varepsilon$, clip radius $R$
\STATE Server initialises $\mathcal{P}^{\mathrm{glob}, (0)} = \mathbf{0}$
\FOR{$t = 1, \dots, T$}
    \FOR{each client $k$ \textbf{in parallel}}
        \STATE Retrieve $\mathcal{P}^{\mathrm{glob}, (t-1)}$; abstract graphs via VCSA ($\phi_{\mathrm{type}}$, Eq.~\ref{eq:edge_mask})
        \STATE Train $\theta_k$ for $E$ local epochs with MGMP (Eqs.~\ref{eq:mgmp_local}--\ref{eq:mgmp_fuse}) and loss $\mathcal{L}_{BCE} + \alpha \mathcal{L}_{SCL} + \gamma \mathcal{L}_{\mathrm{sparse}} + \mu \mathcal{L}_{\mathrm{proto}}$ (Eq.~\ref{eq:proto_align})
        \STATE Build clipped prototypes $\mathbf{p}_c^{(k)}$ (Eq.~\ref{eq:prototype}); perturb via Eq.~\eqref{eq:laplace}
        \STATE Upload $\{\hat{\mathbf{p}}_c^{(k)}\}_{c \in \mathcal{C}}$ \COMMENT{$2|\mathcal{C}|d$ floats per round incl.\ download}
    \ENDFOR
    \STATE Server aggregates $\mathcal{P}^{\mathrm{glob}, (t)}$ via Eq.~\eqref{eq:mcfpa} and broadcasts
\ENDFOR
\ENSURE Prototype bank $\mathcal{P}^{\mathrm{glob}, (T)}$; ensemble detector of Eq.~\eqref{eq:ensemble}; spent budget $\varepsilon_{\mathrm{tot}} = T\varepsilon$
\end{algorithmic}
\end{algorithm}

\subsection{Communication Cost}
\label{sec:comm}

Per round, each client uploads and downloads one $(|\mathcal{C}|+1) \times d$ float32 bank: $2 \times 11 \times 128 \times 4\,\mathrm{B} \approx 11\,\mathrm{KB}$ per client per round in our configuration, independent of model size; total server traffic is $K$ times that (e.g., $\approx$220\,KB per round at $K = 20$). Standard FedAvg on the same backbone transmits the full parameter vector ($\sim$1.4\,MB for our deliberately compact GNN; hundreds of MB for code-LM backbones) per client per round.
